\documentclass[11pt]{article}
\usepackage[margin=1in]{geometry}
\usepackage{amsmath,amssymb,amsthm}
\usepackage{graphicx}
\usepackage{xcolor}
\usepackage{microtype}
\usepackage[colorlinks=true,linkcolor=blue!55!black,citecolor=blue!55!black,urlcolor=blue!55!black]{hyperref}

\newtheorem{theorem}{Theorem}
\newcommand{\K}{\mathbb K}
\newcommand{\N}{\mathbb N}
\newcommand{\cC}{\mathfrak c}

\title{A reflexive counterexample to the outside-$c_0$ spaceability question}
\author{Solution packet for arXiv:2406.06859v2}
\date{12 August 2026}

\begin{document}
\maketitle

\noindent\textbf{Status.}
Candidate full counterexample, likely valid, awaiting human review.

\section{Source question and its interpretation}

For a sequence space $F\subseteq\K^{\N}$, Alves and Ribeiro write
\cite{AlvesRibeiro}
\[
 Z(F)=\{x\in F:x\text{ has only finitely many zero coordinates}\}.
\]
They prove, under a hypothesis involving $c_0$, that
$F\setminus Z(F)$ is $(\alpha,\cC)$-spaceable if and only if $\alpha$ is
finite, and end with the following question on source PDF page 20.

\begin{center}
 \includegraphics[width=0.98\textwidth]{figures/open_problem_crop.png}
\end{center}

Literally, no vector subspace can be contained in $\ell_\infty\setminus c_0$,
because every vector subspace contains $0\in c_0$.  The standard coherent
reading is
\begin{equation}\label{eq:outside}
 F\cap c_0=\{0\},
 \qquad\text{equivalently}\qquad
 F\setminus\{0\}\subseteq\ell_\infty\setminus c_0.
\end{equation}
One might instead read the question as merely requiring $c_0\nsubseteq F$, or
even as requiring that $F$ contain no subspace isomorphic to $c_0$.  The same
counterexample satisfies all three interpretations.

Recall that a set $A$ in a topological vector space is
$(\alpha,\beta)$-spaceable if it is $\alpha$-lineable and every
$\alpha$-dimensional subspace $E\subseteq A\cup\{0\}$ is contained in a
closed $\beta$-dimensional subspace $G\subseteq A\cup\{0\}$.

\section{Proof intuition}

Embed $\ell_2$ isometrically into $\ell_\infty$ using a countable norming
family of functionals.  Instead of recording each functional once, repeat its
value on an infinite coordinate block, and reserve one further infinite block
on which every image is zero.  A nonzero image cannot tend to zero because at
least one nonzero functional value repeats infinitely often, so the range
meets $c_0$ only at zero.  Yet the reserved block gives every image infinitely
many zero coordinates.  Thus $Z(F)$ is empty and the exotic set
$F\setminus Z(F)$ is simply the entire closed space $F$.

\section{The counterexample}

\begin{theorem}\label{thm:counterexample}
Over $\K=\mathbb R$ or $\mathbb C$, there exists a closed subspace
$F\subseteq\ell_\infty$ such that:
\begin{enumerate}
 \item $F$ is linearly isometric to $\ell_2$, hence is reflexive and contains
       no subspace isomorphic to $c_0$;
 \item $F\cap c_0=\{0\}$;
 \item $Z(F)=\varnothing$.
\end{enumerate}
Consequently, for every cardinal $1\le\alpha\le\cC$, the set
$F\setminus Z(F)$ is $(\alpha,\cC)$-spaceable.  In particular the
finite-$\alpha$ characterization proposed in the source question does not
extend to subspaces outside $c_0$.
\end{theorem}

\begin{proof}
Let $H=\ell_2$ and choose a sequence $(\varphi_k)_{k\ge1}$ in the unit sphere
of $H^*$ which is norming:
\begin{equation}\label{eq:norming}
 \sup_{k\ge1}|\varphi_k(x)|=\|x\|_2\qquad(x\in H).
\end{equation}
For example, identify $H^*$ with $H$ and take the functionals induced by a
countable dense subset of the unit sphere.  Density and Cauchy--Schwarz give
\eqref{eq:norming}.

Partition $\N$ into pairwise disjoint infinite sets
\[
 \N=B\,\dot\cup\,A_1\,\dot\cup\,A_2\,\dot\cup\cdots.
\]
One explicit choice is
\[
 B=\{2m-1:m\ge1\},\qquad
 A_k=\{2^k(2m-1):m\ge1\}\quad(k\ge1).
\]
Define $J:H\to\ell_\infty$ coordinatewise by
\begin{equation}\label{eq:J}
 (Jx)_n=
 \begin{cases}
  0,&n\in B,\\
  \varphi_k(x),&n\in A_k.
 \end{cases}
\end{equation}
This map is linear, and \eqref{eq:norming} gives
\[
 \|Jx\|_\infty=\sup_{k\ge1}|\varphi_k(x)|=\|x\|_2.
\]
Thus $J$ is an isometry.  Its range $F=J(H)$ is closed in $\ell_\infty$.

Suppose $Jx\in c_0$.  For each fixed $k$, the coordinates of $Jx$ are equal
to the constant $\varphi_k(x)$ along the infinite, hence unbounded, set
$A_k$.  Convergence to zero forces $\varphi_k(x)=0$ for every $k$.
Equation~\eqref{eq:norming} then gives $x=0$.  Therefore
$F\cap c_0=\{0\}$.

On the other hand, \eqref{eq:J} makes every $Jx$ vanish at every coordinate
in the infinite set $B$.  No member of $F$ has only finitely many zero
coordinates, so $Z(F)=\varnothing$.  Hence
\begin{equation}\label{eq:whole}
 F\setminus Z(F)=F.
\end{equation}

The Hamel dimension of $H$ is $\cC$.  For completeness, $|H|=\cC$ gives the
upper bound, while the vectors
\[
 v_t=(t^n)_{n\ge1}\in\ell_2\qquad(0<t<1)
\]
form a linearly independent family of cardinality $\cC$: every finite
subfamily is independent by the Vandermonde determinant.  Thus
$\dim F=\cC$.

Now fix $1\le\alpha\le\cC$ and any $\alpha$-dimensional subspace
$E\subseteq(F\setminus Z(F))\cup\{0\}=F$.  By \eqref{eq:whole}, the space
$F$ itself is a closed $\cC$-dimensional subspace satisfying
\[
 E\subseteq F\subseteq(F\setminus Z(F))\cup\{0\}.
\]
This is exactly $(\alpha,\cC)$-spaceability.

Finally, $F$ is isometric to the Hilbert space $H$, so it is reflexive.
Since $c_0$ is not reflexive and a Banach-space isomorphic embedding has
closed range, $F$ contains no subspace isomorphic to $c_0$.  This also verifies
the strongest interpretation discussed after \eqref{eq:outside}.
\end{proof}

\section{Scope, verification, and novelty}

The theorem supplies a full counterexample rather than only exploiting the
literal defect in the phrase ``subspace of $\ell_\infty\setminus c_0$.''  It
satisfies the repaired condition \eqref{eq:outside}, the weaker condition
$c_0\nsubseteq F$, and the stronger isomorphic-copy exclusion.

A bounded literature check through 12 August 2026 searched the current arXiv
record and v2 full text, the exact title and arXiv id, the exact open-question
wording, citation-oriented results, and close searches involving
$F\cap c_0$, infinitely many zero coordinates, and
$(\alpha,\cC)$-spaceability.  No later answer was found.  This is a bounded
web/arXiv novelty screen, not an exhaustive bibliographic guarantee.

\medskip
\noindent\textbf{Human-review recommendation.}
High priority.  Check \eqref{eq:norming}, closedness of the isometric range,
the two block arguments, and the final direct use of the definition of
$(\alpha,\cC)$-spaceability.

\begin{thebibliography}{1}
\bibitem{AlvesRibeiro}
D.~Alves and G.~Ribeiro,
\emph{On Sequences with at Most a Finite Number of Zero Coordinates},
arXiv:2406.06859v2 [math.FA], 14 June 2024.
\end{thebibliography}

\end{document}
